\section{Joining two classes while preserving unique extensions}
\label{sec:bridge}

We now choose the abelian subalgebras in the reduced HNN construction.
The resulting inclusion has two features needed in the iteration: every
old pure state extends uniquely, and precisely the prescribed two old
representation classes merge. The proof must control arbitrary old
pure states, even though the defining data consist of only two peaking
sequences.

\begin{theorem}[The bridge]\label{na:bridge}
Let $A$ be a unital $C^*$-algebra with a faithful tracial state $\tau$
and no finite-dimensional irreducible representations. For $i=1,2$,
let $(p_n^i)_{n\geq0}$ be a decreasing sequence of projections, with
$p_0^i=1$, peaking to a pure state $f_i$. Suppose that $f_1\not\sim f_2$
and
\[
 \tau(p_n^1)=\tau(p_n^2)=t_n,
 \qquad 1=t_0>t_1>\cdots>0,
 \qquad t_n\longrightarrow0.
\]
Put $D_i=C^*(p_n^i:n\geq0)$ and let
$\theta:D_1\to D_2$ be determined by $\theta(p_n^1)=p_n^2$.
Form the reduced HNN extension $H(A)$ of Proposition~\ref{na:hnn}
using these data and their trace-preserving conditional expectations.
Then every pure state $h$ of $A$ has a unique state extension
$\widetilde h$ to $H(A)$. These extensions are pure, and for all pure
states $h,k$ of $A$,
\[
 \widetilde h\sim\widetilde k
 \quad\Longleftrightarrow\quad
 h\sim k
 \quad\text{or}\quad
 \{[h],[k]\}=\{[f_1],[f_2]\}.
\]
The extension has a faithful tracial state extending $\tau$ and a
faithful trace-preserving conditional expectation onto $A$. If $A$ is
simple, then $H(A)$ is simple.
\end{theorem}

\begin{proof}
Lemma~\ref{na:peaking-expectation} supplies the expectations onto $D_i$.
The differences $p_n^i-p_{n+1}^i$ are nonzero, mutually orthogonal
projections, so each $D_i$ is the algebra of continuous functions on
$\mathbb N_0\cup\{\infty\}$. Matching these differences defines
$\theta$, and the equality of trace values makes it trace preserving.
Thus the reduced HNN construction applies. Its trace and expectation
already give the corresponding assertions in the theorem. The work
is to prove the assertions about pure states.

\medskip\noindent
\emph{Pass to the two-by-two construction.}
Use the notation of Lemma~\ref{na:free-product} and
Proposition~\ref{na:hnn}:
\[
 F=M_2(A),\qquad G=M_2(D_1),\qquad D=D_1\oplus D_2,
 \qquad Z=F*_D G.
\]
Write $E_D^F:F\to D$ for the expectation, $e_i=e_{ii}$ for the
common diagonal projections, and $s=e_{21}^G$. Recall that
$H(A)=e_1Ze_1$, that $Z=M_2(H(A))$ using the matrix units of $F$,
and that $v=e_{12}^F s$ satisfies $vp_n^1v^*=p_n^2$ in the corner.
Put
\[
 P_n^i=p_n^i e_i\in F,\qquad \phi_i(a)=f_i(a_{ii})\quad(a\in F).
\]
Here $p_n^i e_i$ denotes the matrix with entry $p_n^i$ in position
$(i,i)$ and zero elsewhere. The projections $P_n^i$ peak to $\phi_i$:
their zeroth term is $e_i$, so a state taking value one on them is
supported on that corner and restricts there to $f_i$. The states
$\phi_1,\phi_2$ are pure and inequivalent, by matrix amplification of
the GNS representations of $f_1,f_2$.
Also $F$ has no finite-dimensional irreducible representations. Indeed,
any such representation would restrict along $a\mapsto\operatorname{diag}(a,a)$
to a nonzero finite-dimensional representation of $A$, whose image
would supply a finite-dimensional irreducible representation of $A$.

We shall prove unique extension and the exclusion of unintended joins
for $F\subset Z$, and then return to the corner. By
Lemma~\ref{na:unitary-states}, it suffices to use $\phi_i$ as the
representative of its distinguished equivalence class. More explicitly,
conjugating by a unitary of $F$ bijects the sets of state extensions.
We first establish uniqueness for these representatives and for every
pure state outside their two classes; only afterward will we use the
resulting unique lifts.

Consider pure states $h,k$ of $F$ satisfying one of the two tests
\begin{equation}\label{eq:bridge-tests}
 h=k,
 \qquad\text{or}\qquad
 h\not\sim k\ \text{ and }\quad
 \{[h],[k]\}\ne\{[\phi_1],[\phi_2]\}.
\end{equation}
Whenever either state belongs to $[\phi_i]$, take it to be $\phi_i$
itself. Choose decreasing positive contractive excision nets
$x=x_\lambda$ for $h$ and $y=y_\nu$ for $k$ as in
Lemma~\ref{na:excision}; at $\phi_i$, choose the sequence $P_n^i$
by Lemma~\ref{na:projection-excision}.
All limits involving both $x$ and $y$ below are on the product of the
two index sets.

We will show that $x_\lambda W y_\nu\to0$ in norm for every
nonempty word $W$ in \eqref{na:matrix-word-span}.
When $h=k$, this will force all state extensions to agree; for the
inequivalent pairs in \eqref{eq:bridge-tests}, it will keep the lifted
representations inequivalent. The stable unitary will then join the
designated pair.

\medskip\noindent
\emph{A norm estimate for reduced words.}
We claim that
\begin{equation}\label{eq:bridge-word-decay}
 \|x a_0b_1a_1\cdots b_m a_m y\|\longrightarrow0
\end{equation}
for every $m\geq1$, every $b_j\in\{s,s^*\}$, and every $a_j\in F$
with
\begin{equation}\label{eq:bridge-centering}
 E_D^F(a_j)=0\qquad(0<j<m).
\end{equation}
By the spanning description \eqref{na:matrix-word-span}, these words,
together with $F$, span a dense subspace of $Z$.

Let $P_G$ be the orthogonal projection of the reduced free-product
Hilbert $D$-module $K$ onto the summands starting with a factor in
$L_G^\circ$. For $d\in F$, the decomposition
$K=L_F\otimes_D K^{(F)}$ gives
\begin{equation}\label{eq:bridge-module-bound}
 \|dP_G\|\leq\|E_D^F(d^*d)\|^{1/2}.
\end{equation}
Indeed, a vector $\eta\in P_GK$ has the form $1\otimes\eta$
in this decomposition, and
\[
 \|d\otimes\eta\|^2
 =\|\langle\eta,E_D^F(d^*d)\eta\rangle\|
 \leq\|E_D^F(d^*d)\|\,\|\eta\|^2.
\]
The expectation norm here is the norm in $D$.

For $b\in\{s,s^*\}$ write $b=e_{\ell(b)}be_{r(b)}$.
Thus $(\ell(s),r(s))=(2,1)$ and
$(\ell(s^*),r(s^*))=(1,2)$. For a word in
\eqref{eq:bridge-word-decay}, put
\[
 \ell=\ell(b_1),\qquad r=r(b_m),\qquad
 w=b_1a_1\cdots a_{m-1}b_m.
\]
Each letter of this centered alternating core creates a new leg when
applied, from right to left, to a word not starting with $G$. Therefore
\[
 (1-P_G)w(1-P_G)=0,
 \qquad w=P_Gw+(1-P_G)wP_G.
\]
Set $d=xa_0e_\ell$ and $d'=e_ra_my$. We obtain
\begin{equation}\label{eq:bridge-two-ends}
 \|dwd'\|
 \leq \|dP_G\|\,\|w\|\,\|d'\|
       +\|d\|\,\|w\|\,\|P_Gd'\|.
\end{equation}
Using \eqref{eq:bridge-module-bound} and its adjoint version, it follows
that \eqref{eq:bridge-word-decay} holds whenever both
\begin{equation}\label{eq:bridge-endpoints}
 \begin{split}
 \|E_D^F(e_\ell a_0^*x^2a_0e_\ell)\|&\longrightarrow0,\\
 \|E_D^F(e_r a_my^2a_m^*e_r)\|&\longrightarrow0.
 \end{split}
\end{equation}
This sufficient criterion does not require the pair of endpoint
states to satisfy \eqref{eq:bridge-tests}.

\medskip\noindent
\emph{When an endpoint decays.}
In the first line of \eqref{eq:bridge-endpoints}, increase $x^2$ to $x$.
The resulting elements are a decreasing net of positive continuous
functions on the $\ell$th convergent sequence in the spectrum of $D$;
they vanish on the other component. At a finite atom
$Q=P_n^\ell-P_{n+1}^\ell$, their pointwise limit is
\[
 \frac{\tau_F(Qa_0^*r_ha_0Q)}{\tau_F(Q)}=0.
\]
Here $r_h$ is the support projection in $F^{**}$ and $\tau_F$ is the
normalized matrix trace, extended normally to $F^{**}$. The numerator
vanishes by traciality and Lemma~\ref{na:trace-null}; for example
\[
 0\leq\tau_F(r_h a_0Qa_0^*)
 \leq\|a_0Qa_0^*\|\tau_F(r_h)=0.
\]
At the point $\infty$, the limit is
$\phi_\ell(a_0^*r_ha_0)$. It is zero if $h\not\sim\phi_\ell$,
because the normal extension of the $\phi_\ell$ representation
annihilates $r_h$, as in Lemma~\ref{na:cross-excision}.
If $h=\phi_\ell$, the support identity instead gives
\[
 \phi_\ell(a_0^*r_ha_0)=|\phi_\ell(a_0)|^2.
\]
Thus the first endpoint decays unless $h=\phi_\ell$, and it still
decays in that case when $\phi_\ell(a_0)=0$.

The uniform convergence used here is Dini's theorem for monotone nets
\cite[Theorem~1.5, p.~306]{NaimpallyPeters2013}. We recall its compactness
argument. For a fixed
$\varepsilon>0$, the open sets where a chosen function is below
$\varepsilon$ cover the compact spectrum. Choose a finite subcover,
and then an index dominating its finitely many indices. Monotonicity
makes every later function uniformly smaller than $\varepsilon$.
The second line of \eqref{eq:bridge-endpoints} is symmetric: it
decays unless $k=\phi_r$, and also decays in that case when
$\phi_r(a_m)=0$.

\medskip\noindent
\emph{Remove the matched endpoints.}
Call the left endpoint matched when $h=\phi_\ell$, and the right
endpoint matched when $k=\phi_r$. At each matched endpoint write
\[
 a_0=\phi_\ell(a_0)1+a_0^0,
 \qquad\text{or}\qquad
 a_m=\phi_r(a_m)1+a_m^0.
\]
Expand the word using these decompositions. In each resulting term,
an unmatched endpoint already satisfies its estimate in
\eqref{eq:bridge-endpoints}, and a matched endpoint with its
zero-value coefficient does as well. We treat the scalar terms by
moving the corresponding peak through its adjacent letter:
\begin{equation}\label{eq:bridge-transfer}
 P_n^{\ell(b)}b=bP_n^{r(b)},
 \qquad bP_n^{r(b)}=P_n^{\ell(b)}b.
\end{equation}
These identities hold in $G$: the diagonal embedding of $D_2$ uses
$\theta^{-1}$, so the two projections have the same matrix entry
$p_n^1$ there.

In the following displays, an untransferred endpoint coefficient
means the coefficient chosen in the expansion. For $m\ge2$, after
suppressing its scalar coefficient, a left scalar term becomes
\[
 P_n^{\ell(b_1)}b_1a_1\cdots b_ma_my
 =b_1\bigl(P_n^{r(b_1)}a_1b_2\cdots b_ma_my\bigr).
\]
The peak now belongs to the other distinguished state, and $b_1$
is a bounded outside factor. A right scalar choice similarly moves
$P_t^{r(b_m)}$ to $P_t^{\ell(b_m)}$ and leaves $b_m$ as a bounded
factor on the right. When both ends are transferred and $m\ge3$,
the resulting expression,
again up to its scalar coefficient, is
\[
 b_1\bigl(P_n^{r(b_1)}a_1b_2\cdots
 b_{m-1}a_{m-1}P_t^{\ell(b_m)}\bigr)b_m.
\]
Suppose at least one core letter remains after the transfers. Each
newly exposed endpoint coefficient is one of the original internal
$a_j$. By \eqref{eq:bridge-centering},
$\phi_1(a_j)=\phi_2(a_j)=0$: these states evaluate the two components
of $E_D^F(a_j)$ at $\infty$. Thus every new endpoint satisfies its
estimate, whether or not its transferred state matches the adjacent
letter. Endpoints that were not transferred satisfy their estimates
by the preceding expansion. Applying \eqref{eq:bridge-endpoints}
to the remaining word therefore proves decay of each such term.
This covers no transfers, one transfer when $m\ge2$, and two
transfers when $m\ge3$, including a remaining core of one letter.
The sufficient endpoint criterion requires no assumption that the
transferred pair of states satisfies \eqref{eq:bridge-tests}.

It remains to check the cases in which the core disappears.
If $m=1$, both ends cannot match: the range and source of $s$ or
$s^*$ are different, so double matching would be the designated pair
excluded in \eqref{eq:bridge-tests}. For a single left transfer,
say $h=\phi_i$, the remaining compression, up to a scalar and a
bounded outside letter, is
\[
 P_n^j a_1y_\nu,\qquad j\ne i.
\]
The state $k$ is inequivalent to $\phi_j$. In the equal-state test it
equals $\phi_i$; in the other test, equivalence to $\phi_j$ would
again be the excluded designated pair. Hence this expression tends
to zero by Lemma~\ref{na:cross-excision}. A single right transfer
is symmetric.

Finally let $m=2$ and transfer at both ends. Both original endpoints
matched, so \eqref{eq:bridge-tests} forces $h=k=\phi_i$ for some
$i$. Since each letter switches the two matrix indices,
\[
 \ell(b_1)=r(b_2)=i,
 \qquad r(b_1)=\ell(b_2)=j\ne i.
\]
The remaining compression is $P_n^j a_1P_t^j$, with
$\phi_j(a_1)=0$. If $q=\min(n,t)$, nestedness and excision give
\[
 \|P_n^j a_1P_t^j\|
 \leq\|P_q^j a_1P_q^j\|\longrightarrow0.
\]
This is convergence on the product of the two sequence indices.
Together with cross-excision, whose conclusion also holds on the
product net, these cases prove \eqref{eq:bridge-word-decay} in full.

\medskip\noindent
\emph{Unique extensions and the exclusion of other joins.}
In the equal-state test, let $\psi$ be any state of $Z$ extending $h$.
Its GNS vector is fixed by every $x_\lambda$ and $y_\nu$: a positive
contraction $x$ of state value one satisfies
$\|(1-\pi_\psi(x))\xi_\psi\|^2\leq\psi(1-x)=0$.
Consequently, for every tested word $W$,
\[
 \psi(W)=\psi(x_\lambda Wy_\nu)\longrightarrow0.
\]
Its values on $F$ are already prescribed. The dense span above
therefore proves uniqueness; existence also follows directly by
composing $h$ with the expectation onto $F$. Conjugacy by unitaries
in $F$ now establishes uniqueness for every old pure state, including
all states in the two distinguished classes. Their unique extensions
are pure by Lemma~\ref{na:unique-lift}.

Next take an inequivalent pair allowed in
\eqref{eq:bridge-tests}. If its unique lifts were equivalent, realize
them by unit vectors $\xi_h,\xi_k$ in one irreducible representation
$\pi$ of $Z$. The represented excision elements fix the corresponding
vectors, so
\[
 \langle\pi(z)\xi_k,\xi_h\rangle
 =\langle\pi(x_\lambda zy_\nu)\xi_k,\xi_h\rangle.
\]
For $z\in F$, the right side tends to zero by cross-excision; for a
tested nonempty word it tends to zero by
\eqref{eq:bridge-word-decay}. Thus the matrix coefficient vanishes on
all of $Z$. This contradicts cyclicity of $\xi_k$ in the irreducible
representation. Unitary conjugacy within the old classes extends this
exclusion to all representatives. Equivalent old states, on the other
hand, have equivalent unique lifts by
Lemmas~\ref{na:unitary-states} and~\ref{na:unique-lift}.

\medskip\noindent
\emph{Return to the corner and join the designated pair.}
States of $A=e_1Fe_1$ identify with states of $F$ taking value one on
$e_1$, by composition with compression. The same correspondence holds
between $H(A)=e_1Ze_1$ and $Z$, and it commutes with restriction.
Every extension to $Z$ of an $e_1$-supported state of $F$ is still
$e_1$-supported. Hence this correspondence bijects extension sets and
transfers unique extension to $A\subset H(A)$.

Under $F=M_2(A)$ and $Z=M_2(H(A))$, these supported states have the
matrix amplifications of their corner GNS representations, with the
representing vector in the first summand. Purity and equivalence of
classes are therefore preserved. The matrix-flip unitary of $F$
moves the representative $\phi_2$ from the second summand to the
first, so its class corresponds to $[f_2]$, just as the class of
$\phi_1$ corresponds to $[f_1]$. The exclusion of all other joins
thus descends to the corner.

To see that the designated join actually occurs, use the stable
unitary relation. The state $\widetilde f_2\circ\Ad v$ takes value
one on every $p_n^1$, so its restriction to $A$ is $f_1$. Unique
extension, now proved, gives
\[
 \widetilde f_1=\widetilde f_2\circ\Ad v.
\]
This proves the asserted equivalence criterion. If $A$ is simple,
Lemma~\ref{na:simple-extension} applies to the faithful expectation
and the unique-extension property, proving that $H(A)$ is simple.
\end{proof}
